% Paper-ready replacement Results section for the outcome-blind rigorous-v2 campaign.
% Source run: GitHub Actions 31313459321, artifact rigorous-v2_primary-analysis.
% Frozen code/profile commit: 60064d92739ccff423e085fc8a5f8c6489e764e6.
% This section uses only rigorous-v2 results; legacy numerical values are not inputs.

\section{Results of the outcome-blind primary campaign}
\label{sec:results-v2}

The prespecified primary campaign was executed at depth $d=6n$ for
$n=6,8,10,12$.  It contained five nonconserving ansatz families, with
16 independently drawn circuits per family and size, and a half-filled
$U(1)$-conserving family with 20 independent circuits per size.  Each
fixed circuit was probed with $M=128$ normalized Gaussian parameter-space
directions.  Thus the primary data set contains 400 independent circuit
instances and 51,200 analytic tangent directions.  The fixed circuit,
not the tangent direction, is the independent statistical unit.

All prespecified technical-validity checks passed.  Every scheduled job
and readout row was present.  The maximum state-normalization error was
$2.52\times10^{-14}$, the maximum horizontal overlap was
$4.51\times10^{-15}$, the maximum probability-normalization error was
$2.51\times10^{-14}$, and the largest absolute probability-tangent sum
was $1.80\times10^{-16}$.  The classical Fisher information remained
below the QFI in every tested direction, and the empirical covariance--
projector deviation bound was satisfied with minimum slack $0.1608$.
No circuit was excluded on the basis of a scientific outcome.

\begin{table*}[t]
\caption{Outcome-blind primary results at $d=6n$.  Generic entries are
equal-weight averages of the five nonconserving ansatz-family means;
circuit instances are resampled within family.  The $U(1)$ entries use
20 independent circuits per size.  Brackets denote 95\% circuit-bootstrap
confidence intervals.  The quantity $d_{\rm eff}=1/\widehat{\mathrm{Tr}(C^2)}$
is reported as a reciprocal diagnostic, not as an unbiased estimator.}
\label{tab:v2-primary}
\centering
\begin{tabular}{c c c c c c}
\hline\hline
$n$ & group & $\rho_1$ & $\rho_2$ & $d_{\rm eff}$ & $F_{\rm full}/F_Q$ \\
\hline
6  & generic & $0.960\,[0.947,0.973]$ & $0.978\,[0.972,0.984]$ & $46.98\,[46.68,47.29]$ & $0.4880\,[0.4865,0.4895]$ \\
6  & $U(1)$  & $2.058\,[2.023,2.095]$ & $1.192\,[1.172,1.211]$ & $9.42\,[9.15,9.68]$ & $0.4663\,[0.4590,0.4739]$ \\
8  & generic & $0.947\,[0.936,0.959]$ & $0.974\,[0.970,0.979]$ & $145.00\,[144.20,145.81]$ & $0.4901\,[0.4893,0.4910]$ \\
8  & $U(1)$  & $4.219\,[4.116,4.318]$ & $1.898\,[1.869,1.927]$ & $17.05\,[16.66,17.42]$ & $0.4761\,[0.4687,0.4839]$ \\
10 & generic & $0.944\,[0.934,0.954]$ & $0.967\,[0.964,0.971]$ & $367.50\,[365.45,369.47]$ & $0.4921\,[0.4917,0.4925]$ \\
10 & $U(1)$  & $9.886\,[9.714,10.061]$ & $3.606\,[3.566,3.647]$ & $27.27\,[26.87,27.71]$ & $0.4691\,[0.4648,0.4733]$ \\
12 & generic & $0.943\,[0.934,0.951]$ & $0.963\,[0.960,0.966]$ & $723.72\,[720.27,727.03]$ & $0.4936\,[0.4934,0.4938]$ \\
12 & $U(1)$  & $25.486\,[24.996,25.984]$ & $7.882\,[7.809,7.957]$ & $38.55\,[37.91,39.18]$ & $0.4748\,[0.4723,0.4772]$ \\
\hline\hline
\end{tabular}
\end{table*}

\subsection{Generic low-weight readouts satisfy the prespecified rank-law equivalence criterion}

The rank law was tested as an equivalence statement rather than by a
failure-to-reject argument.  Before observing rigorous-v2 outcomes, the
practical equivalence band was fixed to
$0.90\leq\rho_k\leq1.10$.  For the equal-weight generic-family aggregate,
the full 95\% confidence interval lies inside this band for both one- and
two-body readouts at every tested size.  Specifically,
$\rho_1=0.960,0.947,0.944,0.943$ and
$\rho_2=0.978,0.974,0.967,0.963$ for $n=6,8,10,12$, respectively.
The corresponding leave-one-family-out point estimates shift by at most
$0.0143$, so the family-balanced aggregate is not controlled by a single
ansatz family.

The family-resolved results nevertheless rule out a stronger statement
that every nonconserving architecture is individually equivalent at the
same precision.  The one-body $R_Y$--$R_Z$--CZ-line family fails the
prespecified equivalence criterion at $n=8,10,12$; its means are
$0.900$, $0.891$, and $0.886$, respectively, with the latter two clearly
below the lower equivalence boundary.  The Haar-$U(4)$ brickwork family
is also inconclusive at $n=8$ because its 95\% interval extends slightly
below $0.90$.  These exceptions are retained rather than absorbed into
the pooled statement.  Thus the supported conclusion is a
family-balanced generic rank-typicality statement, not architecture-
universal exact rank control.

\subsection{The symmetry-constrained family exhibits a rapidly growing low-weight excess}

The $U(1)$ family lies far outside the rank-law band and its enhancement
grows strongly over the tested finite-size range.  The one-body
quantity increases from $\rho_1=2.058$ at $n=6$ to
$\rho_1=25.486$ at $n=12$, while the two-body quantity increases from
$\rho_2=1.192$ to $\rho_2=7.882$.  The corresponding physical retention
fractions at $n=12$ are $0.3037\,[0.2980,0.3097]$ and
$0.5551\,[0.5498,0.5603]$ for one- and two-body readouts, compared with
$0.002762\,[0.002738,0.002787]$ and
$0.018344\,[0.018289,0.018400]$ for the family-balanced generic ensemble.

The prespecified $U(1)$-minus-generic contrasts are therefore large and
well resolved.  At $n=12$, the one-body retention difference is
$0.30097\,[0.29517,0.30700]$ and the two-body difference is
$0.53671\,[0.53165,0.54186]$.  The corresponding enhancement differences
are $24.543\,[24.058,25.038]$ and $6.919\,[6.847,6.992]$.
All of these confidence intervals exclude zero without any post-hoc
metric selection.

\subsection{Rank typicality does not imply global isotropy}

The second spectral moment shows that the generic visible-score
covariance becomes increasingly nonflat even while its low-weight
readouts satisfy the rank-law equivalence criterion.  For the
family-balanced generic ensemble,
$N\,\widehat{\mathrm{Tr}(C^2)}$ increases from
$1.350\,[1.340,1.359]$ at $n=6$ to
$6.165\,[6.116,6.216]$ at $n=12$, whereas
$d_{\rm eff}/N$ falls from $0.746\,[0.741,0.751]$ to
$0.1767\,[0.1759,0.1776]$.  Hence the covariance is not approaching the
flat isotropic reference in this finite-size window.

The $U(1)$ covariance is substantially more concentrated.  Its normalized
second moment grows from
$N\,\widehat{\mathrm{Tr}(C^2)}=2.026\,[1.968,2.085]$ at $n=6$ to
$23.975\,[23.587,24.380]$ at $n=12$, while
$d_{\rm eff}/N$ decreases from $0.496\,[0.481,0.510]$ to
$0.0418\,[0.0411,0.0425]$.  These results directly separate low-weight
rank typicality from global covariance isotropy: the generic ensemble
can satisfy the prescribed low-rank overlap law even when its score
covariance is strongly anisotropic, whereas the symmetry-constrained
family is both more anisotropic and far more visible to low-weight
readouts.

\subsection{The separation is primarily a readout-orientation effect}

The complete computational-basis record retains an order-one fraction of
the QFI in both ensembles.  Across $n=6,8,10,12$,
$F_{\rm full}/F_Q$ ranges from $0.4880$ to $0.4936$ in the generic
ensemble and from $0.4663$ to $0.4761$ in the $U(1)$ ensemble.  The
$U(1)$ values are statistically lower by approximately $0.014$--$0.023$
in the prespecified contrasts, but this modest full-record difference is
far too small to account for the much larger separation after projection
onto one- and two-body readouts.  The dominant effect is therefore the
relative geometry between the visible tangent covariance and the
physical low-weight feature subspace.

\subsection{Inference status}

Every primary family/size/readout cell passes the prespecified independent-
circuit count and precision flags for the principal quantities in
Table~\ref{tab:v2-primary}.  These results are therefore the primary
outcome-blind finite-size estimates from the frozen campaign.  The
protocol nevertheless requires an independent tangent-sample convergence
study before $M=128$ endpoints are promoted to the strongest final
claims.  Consequently, the numerical statements above are reportable as
primary results, while the final wording of convergence-sensitive claims
must remain conditional until the separately frozen $M=32,64,128,256$
campaign is completed.  No result in the present size range is interpreted
as an asymptotic proof.
